\subsection{Study 3: Combining Spatial and Temporal Information Enhances Understanding of Ant Foraging Behavior}

This study utilized a series of time-ordered images to create heatmaps illustrating the spatial distribution of ants over time (Figure \ref{fig:s3_heatmap_small}). In subsets \textit{B01} and \textit{B02}, three separate trials (one from B01 and two from B02) were conducted to assess the impact of liquid food dispensing design on ant foraging behavior (Figure \ref{fig:s3_heatmap_small}a). In the first trial, a triangular filter paper was placed to distribute sugar solution from a 5 mL vial. In the second and third trials, a thread was used to dispense sucrose solutions from a 15 mL Falcon tube. The heatmap for the first trial indicates that ants were evenly distributed around the filter paper, with higher densities observed along the edges between the filter paper and the petri dish. For the thread-guided strategy, both trials showed high activity around the thread and the cap of the tube, with the second trial displaying a more concentrated ant distribution along the thread. Analysis of ant distribution over time indicated higher numbers of foraging ants in the first trial, suggesting that the triangular filter paper, which provides a larger total area for liquid food dispensing compared to threads used in other trials, may more effectively attract ants.

Another experiment using subset \textit{B03} explored macronutrient preference (i.e., sucrose (labeled 'S') and peptone (labeled 'P')) in OHAV-1 infected \emph{T. sessile} (Figure \ref{fig:s3_heatmap_small}b). Both infected and uninfected ants were observed, and their preferences were analyzed by examining their distribution relative to a hypothetical line drawn between two food sources (petri dishes). The heatmap suggested a potential preference of infected ants for sucrose, whereas uninfected ants showed no clear preference between the two macronutrients. Time-series analysis of ant presence indicated that infected ants exhibited increased foraging activity primarily within the first five hours, followed by a notable decline. These observations demonstrate the CV system's capability to simultaneously capture spatial and temporal patterns, offering valuable insights into ant foraging behaviors that might be challenging to document through manual counting alone, especially in larger ant colonies. However, further statistical analysis is needed to draw robust conclusions.

The dense ant images from experiment subset \textit{Dense} were visualized in a heatmap (Figure \ref{fig:s3_heatmap_dense}). This experiment examined the impact of SINV-3 infection on ant foraging behaviors, but with a significantly denser ant population and an ant nest cell positioned at the center of the image. The heatmap does not show a strong difference in foraging intensity between uninfected and infected ants; both groups display hotspots in the four corners of the container, reflecting their natural tendency to explore the environment. Additionally, the heatmap indicates that worker ants are more concentrated around water sources compared to food sources and the nest area. This information is essential for understanding ant foraging behavior and can be applied to optimize bait placement strategies in pest control applications.


\begin{figure}[H]
    \centering
    \includegraphics[width=\textwidth]{figure_10.jpg}
    \caption{Heatmap of aggregated ant activity across multiple time points for subsets (a) \textit{B01}, \textit{B02} and (b) \textit{B03}. The top row shows the original images. In panel (b), the letters 'P' and 'S' represent peptone and sucrose, respectively, corresponding to the left and right petri dishes. The middle row shows the corresponding heatmaps illustrating ant activity over time. The color gradient represents the scaled number of ant detections, with white/yellow indicating high activity and black indicating low activity. The bottom row presents line plots of ant presence over time, with the x-axis representing hours and the y-axis representing the number of detected ants.}
    \label{fig:s3_heatmap_small}
\end{figure}

\textit{Alt text}: Visualizations of aggregated ant activity across multiple time points for experiments (a) \textit{B01}, \textit{B02} and (b) \textit{B03}. Top row: Original images. Middle: Heatmaps showing where ants were active; brighter colors indicate more activity. Bottom row: Line graphs of ant counts over time.

\begin{figure}[H]
    \centering
    \includegraphics[width=\textwidth]{figure_11.jpg}
    \caption{Heatmap of ant distribution in subset \textit{Dense}. The left panel displays the uninfected group, while the right panel shows the SINV-3 infected group. The color gradient represents the scaled number of ant detections, with white/yellow indicating high activity and black indicating low activity.}
    \label{fig:s3_heatmap_dense}
\end{figure}

\textit{Alt text}: Heatmaps comparing ant activity in subset \textit{Dense} between the uninfected group (left) and the SINV-3 infected group (right); brighter colors indicate higher ant activity levels.

